\section{Impossibility and Non-Identifiability Results}
\label{sec:impossibility}

This section establishes formal limits of KPI-based observation for the
enterprise continuum \(K_{EA}\) under the Ontology of Continua and the frozen
executable specification \texttt{ETS.K\_EA.v1.1}.
All results are strictly diagnostic and negative in nature: they delimit what
KPI-based observation cannot identify even in principle.

The formal structure follows the classical observation--identification
paradigm: observation is modeled as a map from a state space to an observation
space, and identifiability corresponds to injectivity of that map on a domain of
interest (see, e.g., classical treatments in control and system identification
\cite{kalman1960,ljung1999,walter1997}).
The specific contribution of the OC/ETS framework is to instantiate this
paradigm for enterprise continua with explicit observability constraints.

\subsection{Existence of KPI-Indistinguishable States}

\begin{theorem}[Non-Identifiability in the Absence of Injectivity]
\label{thm:nonident}
Let \(K=(\kappa_1,\ldots,\kappa_m)\) be any finite admissible KPI family with
induced readout map
\[
\pi_K:\Omega_{\mathrm{obs}}\to\Sigma_{\mathrm{obs}}.
\]
If \(\pi_K\) is not injective on \(\Omega_{\mathrm{obs}}\), then there exist
distinct states
\(\omega_1\neq\omega_2\in\Omega_{\mathrm{obs}}\) such that
\[
\pi_K(\omega_1)=\pi_K(\omega_2).
\]
\end{theorem}

\begin{proof}
Non-injectivity of \(\pi_K\) means, by definition, that there exist at least two
distinct elements of \(\Omega_{\mathrm{obs}}\) that are mapped to the same
observation.
These states are therefore indistinguishable under KPI observation.
\end{proof}

Although elementary, this statement is foundational: it makes explicit that
non-identifiability is the default situation whenever injectivity of the
observation map is not guaranteed.

\subsection{Absence of Universal KPI Completeness}

\begin{theorem}[No Universal KPI Completeness]
\label{thm:no-complete}
There exists no finite admissible KPI family \(K\) such that the associated
readout map \(\pi_K\) uniquely identifies all states in
\(\Omega_{\mathrm{obs}}\), unless injectivity is imposed as an additional
constraint on the admissible state space.
In particular, neither OC nor \texttt{ETS.K\_EA.v1.1} implies injectivity of any
finite KPI readout map on \(\Omega_{\mathrm{obs}}\).
\end{theorem}

\begin{proof}[OC/ETS-grounded argument]
The Ontology of Continua and \texttt{ETS.K\_EA.v1.1} fix the structural components
of the enterprise continuum, including its state space, admissible observable
domain, axes, flows, thresholds, cycles, and the factorization of the continuum
measure \(k_{EA}\).
They do not contain any axiom or admissibility condition requiring a user-chosen
KPI readout map \(\pi_K\) to be injective on \(\Omega_{\mathrm{obs}}\).

This mirrors the classical situation in system identification, where
identifiability is not automatic but must be established or enforced through
additional structural assumptions, excitation conditions, or model
restrictions (see, e.g., \cite{ljung1999,walter1997}).
Within OC/ETS, no such enforcing assumption is present.
Hence universal KPI completeness cannot be concluded from OC/ETS alone.
\end{proof}

\subsection{Structural Indistinguishability as an Intrinsic Feature}

\begin{proposition}[Loss of Structural Information Along KPI Fibers]
\label{prop:fiber-loss}
Let \(K\) be any admissible KPI family and let
\(\omega_1,\omega_2\in\Omega_{\mathrm{obs}}\) lie in the same fiber of
\(\pi_K\).
Then all OC/ETS-structural information distinguishing \(\omega_1\) from
\(\omega_2\) is unrecoverable from KPI values alone.
\end{proposition}

\begin{proof}
If \(\omega_1\) and \(\omega_2\) lie in the same fiber, then
\(\pi_K(\omega_1)=\pi_K(\omega_2)\).
Any downstream procedure whose input depends only on KPI values receives
identical input for \(\omega_1\) and \(\omega_2\).
Consequently, all structural variation within a fiber is eliminated once
observation factors through the quotient
\(\Omega_{\mathrm{obs}}/\!\sim_K\).
\end{proof}

This fiber-based loss of information is the direct analogue of classical
projection-induced ambiguity in partially observed systems
(cf.\ indistinguishable states in linear and nonlinear observation theory
\cite{kalman1960,slotine1991}).

\subsection{Independence from Statistical or Algorithmic Enrichment}

\begin{proposition}[Post-Processing Cannot Restore Identifiability]
\label{prop:no-pp}
Let \(g:\Sigma_{\mathrm{obs}}\to Z\) be any deterministic post-processing map,
including aggregation, smoothing, scoring, embedding, or transformation.
Then the composite map \(g\circ\pi_K\) is identifying on a domain
\(D\subseteq\Omega_{\mathrm{obs}}\) only if \(\pi_K\) is identifying on \(D\).
\end{proposition}

\begin{proof}
If \(\pi_K(\omega_1)=\pi_K(\omega_2)\), then
\((g\circ\pi_K)(\omega_1)=(g\circ\pi_K)(\omega_2)\) for any deterministic function
\(g\).
Thus post-processing cannot separate states that are already
KPI-indistinguishable.
Identifiability of \(g\circ\pi_K\) therefore implies identifiability of
\(\pi_K\) on the same domain.
\end{proof}

This reflects a standard principle in identification theory: no deterministic
transformation of insufficient data can recover information that was never
observed in the first place \cite{ljung1999,walter1982}.

\subsection{Observability Versus Identifiability}

\begin{proposition}[Observability Does Not Imply Identifiability]
\label{prop:obs-vs-id}
ETS-level observability guarantees only that KPI readouts are defined on
\(\Omega_{\mathrm{obs}}\).
It does not imply that KPI readouts uniquely identify states in
\(\Omega_{\mathrm{obs}}\).
\end{proposition}

\begin{proof}
Observability, as used here, is an admissibility condition ensuring that readouts
are defined on the domain \(\Omega_{\mathrm{obs}}\).
Identifiability is a distinct property, equivalent to injectivity of the readout
map on a domain of interest.
OC and \texttt{ETS.K\_EA.v1.1} do not equate these notions and do not enforce
injectivity.
This separation mirrors the classical distinction between observability and
identifiability in control and system identification theory
\cite{kalman1960,ljung1999}.
\end{proof}
